
Una vez que la fuente se encuentre en el piso y siempre que se disponga de los accesorios necesarios para la recuperación de la fuente, se podrán iniciar las operaciones.

La operación de recuperación deberá practicarse antes de realizar efectivamente con la fuente, con objeto de adquirir habilidad para ello.

% Steps for source recovery
\subsubsection*{Pasos para la Recuperación de la Fuente}
Una vez que se inicie la recuperación de la fuente, la operación debe ser rápida (en segundos) siguiendo los siguientes pasos:
\begin{enumerate}
    \item Coloque el contenedor para recuperación a una distancia de la fuente igual a dos veces la distancia de la caña a la fuente.
    \item Coloque el embudo sobre el contenedor.
    \item La persona que realice las operaciones deberá colocarse entre el contenedor y la fuente a la distancia que la caña lo permita.
    \item Fije la fuente al imán de la caña.
    \item Haga un giro con la caña, manteniéndola horizontalmente para verificar que se encuentra bien unida, esto evitará que la fuente quede cerca si se saliera de la caña.
    \item Coloque la fuente sobre el embudo con el acople (o gancho) hacia arriba.
    \item Sacuda la caja para que la fuente caiga dentro del contenedor.
    \item Protegiéndose con el contenedor, flexione el cable de la fuente y engánchelo al cable del impulsor (reel). En caso necesario asegure la unión con una capa de cinta adhesiva industrial.
    \item Introduzca la fuente al contenedor, accionando la manguera del impulsor.
    \item Quite la manguera del contenedor original, acérquelo al contenedor de recuperación, pase a través de él el cable del impulsor.
    \item Asegure la fuente en su posición dentro del contenedor original.
\end{enumerate}

% Explanation of the safety measures
\textbf{Nota:} Esta secuencia de pasos asegura que la fuente radiactiva sea manipulada de manera segura y eficiente, minimizando el riesgo de exposición accidental a la radiación.
